\documentclass{beamer}

% --- PREAMBLE ---
\usepackage[utf8]{inputenc}
\usepackage[indonesian]{babel}
\usepackage{graphicx}
\usepackage{listings} % Untuk kode

% Konfigurasi untuk blok kode
\lstset{
  basicstyle=\small\ttfamily,
  breaklines=true,
  backgroundcolor=\color{gray!10},
  frame=single,
  rulecolor=\color{black!40}
}

% --- TEMA PRESENTASI ---
% Anda bisa mengganti tema sesuai selera (contoh: Madrid, Boadilla, Warsaw, etc.)
\usetheme{Madrid}
\usecolortheme{beaver}

% --- INFORMASI JUDUL ---
\title[Coding \& DL untuk Analisis LIBS]{Penggunaan Coding dan Deep Learning untuk Analisis Spektrum LIBS}
\subtitle{Studi Kasus Implementasi Teknis Berbasis Data Sintetis}
\author{Oleh: Birrul Walidain}
\institute{Program Studi Fisika \\ Fakultas Matematika dan Ilmu Pengetahuan Alam \\ Universitas Syiah Kuala}
\date{Agustus 2025}

\begin{document}

% --- SLIDE 1: HALAMAN JUDUL ---
\begin{frame}
  \titlepage
\end{frame}

% --- Menambahkan section untuk membuat agenda otomatis ---
\section*{Pendahuluan}
\section{Pondasi: Ekosistem Python}
\section{Implementasi Coding: Generasi Dataset}
\section{Implementasi Coding: Manajemen Data}
\section{Implementasi Deep Learning: Model \& Pelatihan}
\section{Implementasi Deep Learning: Evaluasi}
\section{Kesimpulan}

% --- SLIDE 2: AGENDA PRESENTASI ---
\begin{frame}
  \frametitle{Agenda Presentasi}
  \tableofcontents
\end{frame}

% --- SLIDE 3: KONTEKS MASALAH ---
\begin{frame}
  \frametitle{Konteks: Tantangan Analisis LIBS}
  
  \begin{block}{Masalah Fundamental}
    Analisis LIBS konvensional terhambat oleh \textbf{kompleksitas spektrum}, \textbf{variabilitas sinyal}, dan ketergantungan pada proses \textbf{kalibrasi} yang mahal dan lambat.
  \end{block}
  
  \begin{alertblock}{Solusi Berbasis Komputasi (Paradigma Baru)}
    Mengembangkan metode analisis \textbf{bebas kalibrasi} dengan melatih model \textit{Deep Learning} pada \textbf{dataset sintetis} skala besar. Pendekatan ini sepenuhnya bergantung pada implementasi \textit{coding} yang efisien.
  \end{alertblock}
\end{frame}

% --- SLIDE 4: PONDASI PYTHON ---
\begin{frame}
  \frametitle{Pondasi: Ekosistem Python untuk Fisika Komputasi}
  
  Pemilihan Python bukan tanpa alasan. Ekosistemnya menyediakan semua alat yang dibutuhkan untuk menjembatani fisika dan \textit{deep learning}.
  
  \begin{itemize}
    \item \textbf{Komputasi Numerik (\texttt{NumPy}):} 
    Menjadi inti dari semua operasi data spektral. Kalkulasi fisika seperti Persamaan Saha-Boltzmann dieksekusi secara efisien menggunakan \textbf{operasi vektor}, menghindari \textit{loop} `for` yang lambat.
    
    \item \textbf{Manajemen Data (\texttt{h5py} \& \texttt{Pandas}):} 
    Digunakan untuk mengelola 50.000 sampel data secara terstruktur. \texttt{h5py} krusial untuk menangani dataset yang lebih besar dari memori RAM.
    
    \item \textbf{\textit{Framework Deep Learning} (\texttt{PyTorch}):} 
    Menyediakan "blok bangunan" modular untuk merancang, melatih, dan mengevaluasi arsitektur jaringan saraf yang kompleks seperti \textit{Informer}.
    
    \item \textbf{Evaluasi \& Visualisasi (\texttt{scikit-learn} \& \texttt{Matplotlib}):} 
    Untuk menghitung metrik performa model secara akurat dan memvisualisasikan hasil untuk interpretasi fisis.
  \end{itemize}
\end{frame}

% --- SLIDE 5: CODING GENERASI DATASET ---
\begin{frame}
  \frametitle{Implementasi Coding: Generasi Dataset Sintetis}
  
  Tahap ini adalah tentang menerjemahkan persamaan fisika ke dalam kode Python yang efisien dan dapat diskalakan.
  
  \begin{block}{Alur Kerja Kode}
    \begin{enumerate}
      \item \textbf{Perencanaan Kombinasi:} Modul standar Python seperti \texttt{itertools} dan logika skrip digunakan untuk menghasilkan "resep" kombinasi elemen yang seimbang secara statistik.
      
      \item \textbf{Kalkulasi Intensitas (Vektorisasi):} Populasi tingkat energi dan intensitas relatif dihitung untuk semua transisi secara bersamaan menggunakan \texttt{NumPy}.
      \begin{lstlisting}[language=Python]
# Contoh konseptual vektorisasi dengan NumPy
intensities = (g_k * A_ki * np.exp(-E_k / (k_B * T))) / Z
      \end{lstlisting}
      
      \item \textbf{Pelebaran Garis (\textit{Broadening}):} Spektrum "batang" yang ideal diubah menjadi spektrum yang realistis dengan melakukan konvolusi menggunakan fungsi dari pustaka \texttt{SciPy} (\texttt{scipy.signal.convolve}).
    \end{enumerate}
  \end{block}
\end{frame}

% --- SLIDE 6: CODING MANAJEMEN DATA ---
\begin{frame}
  \frametitle{Implementasi Coding: Manajemen \& Pemuatan Data}
  
  Dataset yang besar memerlukan strategi penanganan yang efisien untuk proses pelatihan.
  
  \begin{block}{Penyimpanan dengan \texttt{h5py}}
  Seluruh 50.000 spektrum dan labelnya disimpan dalam satu file HDF5. Ini memungkinkan \textbf{akses data yang sangat cepat} dan pemuatan data secara bertahap (\textit{on-the-fly loading}) tanpa membebani memori.
  \end{block}
  
  \begin{block}{Pemuat Data Kustom dengan \texttt{PyTorch}}
  Kelas \texttt{Dataset} kustom diimplementasikan di PyTorch untuk menjadi jembatan antara file HDF5 di disk dan model selama pelatihan.
  \begin{lstlisting}[language=Python]
from torch.utils.data import Dataset
import h5py

class LIBSDataset(Dataset):
    def __init__(self, h5_path):
        self.h5_file = h5py.File(h5_path, 'r')
        self.spectra = self.h5_file['spectra']
        self.labels = self.h5_file['labels']

    def __len__(self):
        return len(self.spectra)

    def __getitem__(self, idx):
        # Mengambil satu sampel data
        spectrum = self.spectra[idx]
        label = self.labels[idx]
        return spectrum, label
  \end{lstlisting}
  \end{block}
\end{frame}

% --- SLIDE 7: IMPLEMENTASI DEEP LEARNING ---
\begin{frame}
  \frametitle{Implementasi Deep Learning: Arsitektur Model \& Pelatihan}
  
  \begin{block}{Arsitektur Model Informer dengan \texttt{PyTorch}}
    \begin{itemize}
        \item Model dibangun secara modular menggunakan kelas \texttt{nn.Module} dari PyTorch.
        \item Komponen utamanya adalah \textbf{mekanisme atensi (\textit{attention})} yang dioptimalkan untuk sekuens panjang. Efisiensi komputasi (kompleksitas $O(L \log L)$) lebih baik dibandingkan Transformer standar ($O(L^2)$).
        \item Tujuannya agar model dapat menghubungkan puncak-puncak spektral yang letaknya berjauhan secara efisien.
    \end{itemize}
  \end{block}

  \begin{block}{Teknik Pelatihan}
    \begin{itemize}
      \item \textbf{\textit{Loss Function}:} \textbf{Focal Loss} digunakan agar model lebih fokus belajar dari sinyal-sinyal elemen yang lemah atau sulit terdeteksi.
      \item \textbf{Optimizer:} \textbf{AdamW} dipilih sebagai algoritma optimasi modern yang membantu model belajar lebih cepat dan stabil.
      \item \textbf{Akselerasi Pelatihan:} Proses dipercepat signifikan menggunakan \textbf{\textit{Automatic Mixed Precision}} (\texttt{torch.amp}) untuk memanfaatkan \textit{Tensor Cores} pada GPU modern.
    \end{itemize}
  \end{block}
\end{frame}

% --- SLIDE 8: EVALUASI MODEL ---
\begin{frame}
  \frametitle{Implementasi Deep Learning: Evaluasi Kinerja}

  \begin{block}{Evaluasi Kuantitatif dengan \texttt{scikit-learn}}
    \begin{itemize}
        \item Kinerja model pada data uji diukur menggunakan fungsi \texttt{classification\_report}.
        \item Fungsi ini secara otomatis menghitung metrik-metrik penting seperti \textbf{Precision}, \textbf{Recall}, dan \textbf{F1-Score} untuk setiap elemen, memberikan gambaran lengkap tentang performa model.
    \end{itemize}
  \end{block}
  
  \begin{block}{Pemrosesan Lanjutan: Pencarian Ambang Batas}
    \begin{itemize}
        \item Output model adalah probabilitas (0-1). Sebuah skrip Python dibuat untuk mencari \textbf{ambang batas (threshold)} optimal pada data validasi.
        \item Langkah ini krusial untuk mengubah probabilitas menjadi keputusan akhir (elemen "ada" atau "tidak ada") dan memaksimalkan F1-Score.
    \end{itemize}
  \end{block}
\end{frame}

% --- SLIDE 9: KESIMPULAN ---
\begin{frame}
  \frametitle{Kesimpulan: Peran Sentral Coding \& Deep Learning}
  
  \begin{enumerate}
    \item \textbf{\textit{Coding} sebagai Jembatan:} Kode Python, khususnya dengan \texttt{NumPy} dan \texttt{SciPy}, berhasil menjadi jembatan antara teori fisika plasma dan data latih yang dapat digunakan oleh mesin.
    
    \item \textbf{\textit{Deep Learning} sebagai "Analis" Otomatis:} Arsitektur Informer yang diimplementasikan dengan \texttt{PyTorch} terbukti mampu "belajar" fisika fundamental dari spektrum emisi untuk melakukan identifikasi multi-elemen.
    
    \item \textbf{Ekosistem Terintegrasi:} Keberhasilan pendekatan ini sangat bergantung pada bagaimana ekosistem Python memungkinkan seluruh alur kerja—mulai dari simulasi, manajemen data, hingga pelatihan model AI—dapat berjalan secara mulus dan otomatis.
  \end{enumerate}
  
  \vspace{1em}
  \centering
  \structure{Pendekatan ini membuka jalan bagi metode analisis unsur yang lebih cepat, efisien, dan otomatis.}
\end{frame}

% --- SLIDE 10: TERIMA KASIH ---
\begin{frame}
  \frametitle{Terima Kasih \& Diskusi}
  
  \vfill
  \centering
  \Huge{\textbf{Terima Kasih}}
  \vspace{2em}
  
  \Large{Pertanyaan \& Diskusi}
  \vfill
\end{frame}

\end{document}
